\chapter{The Completion Door, Opened Without Analysis}
\label{chap:completion-door-opened}

Since the sixteenth chapter the construction has carried a door it would not walk through: the
completion, the passage from its finite, refinable tower to the complete space the tower approaches. It
flagged the door honestly each time it leaned on it, calling it the one place the construction seemed to
reach for something it could not itself build. This chapter opens the door, and the news is that opening
it imports nothing. The completion is assembled from the construction's own materials --- the energy norm
it has had since the square-root door, and the separable skeleton the countable ledger of the last
chapter supplies --- without bringing in the analysis a completion is usually thought to require. And once
the door is open, every finite result the construction proved walks through it intact: the split, the
orthogonality, the radiation pinned to the boundary, all lift to the continuum as theorems there. What
the door does not build, it still only grants what the whole book has granted --- that the limits are
reached --- and nothing analytic besides.

\section{The door is built, not imported}
\label{se:the-door-is-built-not-imported}

The completion door the construction opens is a Sobolev--Hilbert door, and its defining furniture is a
norm. The norm is not borrowed: it is the square root of the energy, the very measure the construction
built when it opened the square-root door, satisfying the norm laws the earlier chapters proved --- zero
length exactly at the zero element, and the rest. The finite Hilbert door was already in hand from two
earlier episodes, a normed energy space with everything a pre-Hilbert structure needs; what this chapter
does is extend that finite door to its completion, the Sobolev--Hilbert space the finite one sits densely
inside.

The phrase \emph{without importing analysis} is the load-bearing one, and it is worth being exact about
what it denies. The usual way to complete a normed space is to invoke the completed real line and its
analytic apparatus --- limits, suprema, the whole edifice of the continuum --- as machinery brought in
from outside. The construction refuses that import. It does not reach for a pre-built continuum and graft
its space onto it; it completes its own space by its own means, using only structure it has already
constructed. The completion that results is genuine --- a complete space with the finite one dense in it
--- but it is the construction's completion, assembled from the construction's parts, and the analytic
theory of the reals is nowhere invoked. This is the resolution of the long-flagged door: it was never a
hole through which all of analysis poured in, only a passage the construction had not yet shown it could
build with its own hands.

What this changes about the book's accounting is larger than it first appears. For eight chapters the
completion door has sat in the assumed column with an asterisk, the one place a reader skeptical of
imported infinity might plant a flag and say \emph{here is where you smuggled the continuum in}. Opening
the door without analysis removes the flag. There is no smuggled continuum, because there is no import at
all; there is a separable space the construction built and a completion of it the construction performs,
and a skeptic who accepts the finite normed space and the countable skeleton --- both constructed in
front of them --- has already accepted everything the completion uses but the one grant. The door, in
other words, was never the soft spot it looked like. It looked like the place analysis entered because the
construction had not yet shown the completion could be done from below; now that it has, the door is as
hard as anything else in the book.

\section{The separable skeleton}
\label{se:the-separable-skeleton}

What makes the completion buildable by the construction's own means is the fact the last chapter proved:
the ledger is countable. A pre-Hilbert space with a countable dense skeleton is called separable, and
separability is exactly the property that lets a space be completed by counting rather than by import. The
recipe is elementary. Take Cauchy sequences drawn from the countable skeleton --- sequences whose terms
eventually agree to any tolerance --- and identify two of them whenever they converge together; the
resulting classes are the points of the completion, and the finite skeleton sits inside as the sequences
that are eventually constant. Every step of this uses only countable data and finite tolerances, the
construction's native materials. The countability proved in the previous chapter was not a curiosity; it
was precisely the hinge the completion turns on, the reason the separable skeleton exists and the reason
the Cauchy sequences can be enumerated and managed without an uncountable choice.

It is worth walking the recipe slowly, because its modesty is the whole point. A Cauchy sequence is a list
of skeleton points whose terms eventually huddle within any tolerance one names --- past some stage, all
of them agree to that precision. Two such sequences are deemed to name the same limit point when their
terms eventually huddle together, each near the other. A point of the completion is then nothing more than
a class of such sequences, all converging to the same place; and the original skeleton sits inside as the
sequences that simply stand still, each constant at one of its own points. Every object in this
construction is a sequence of countable data and a tolerance that is a natural number --- nothing
uncountable is ever formed, nothing chosen all at once from an unlistable set. The completion is large,
even uncountable as a whole, but it is assembled entirely from countable pieces, the way the real line
itself can be built from rational Cauchy sequences without presupposing the reals. That is the sense in
which the construction completes its space without importing one: it does to its separable skeleton
exactly what the rationals' completion does to the rationals, and asks no more.

The construction makes the embedding exact: the finite residue maps into the completion, each finite
object carried to the point of the completed space it represents, faithfully and without collision. So
the finite tower is not discarded when the door opens; it is funged into the completion, pooled with the
limit points it approaches into one complete space where the finite and the limiting sit together. This is
the Hilbert completion the experiment under Hilbert's name describes --- a complete inner-product space
reached as the closure of a separable pre-Hilbert skeleton --- built here from below, out of the
countable ledger, rather than postulated from above. The space the construction has been approaching
since it first refined a sample is now a place it can stand in, and it got there by counting.

Separability is a quiet word for a decisive property, and it is worth marking how much rides on the
construction having it for free. A space without a countable dense skeleton cannot be completed by the
sequence recipe at all; its completion, if one wants it, must be summoned by exactly the kind of
non-constructive principle the book has spent twenty-three chapters avoiding. The construction is spared
that entirely, and not by luck: the separability is the countability of the forcing ledger, which the last
chapter proved, which was itself the faithful coding of finite conditions into the naturals. So a chain
runs from the very first finite samples to the completed Sobolev space --- samples are finite conditions,
conditions are countable, countable means separable, separable means completable from below --- and every
link is built, none imported. The completion stands at the end of that chain as its natural terminus, the
place the finite cuts were always heading, reachable because every step toward it was finite.

\section{The finite structure lifts}
\label{se:the-finite-structure-lifts}

An open door is only worth opening if things pass through it, and what passes through this one is every
structural fact the construction earned in the finite setting. The split lifts: the orthogonality of the
geometric and gauge sectors, proved on the finite certificate, holds in the completion, the two sectors
perpendicular in the continuum as they were in the finite space. The interior obstruction lifts: there is
no continuum interior mixed scalar, just as there was none finitely, so the coupling is still forbidden
the bulk. The boundary trace agrees with the terminal residual: the reading taken at the boundary of the
completed space equals the limit of the finite tower's residuals, so the continuum boundary quantity is
the honest limit of the finite ones and not a new thing smuggled in. The residual is Cauchy in the
Sobolev norm, converging in the completed space the way it converged finitely. And the radiation theorem
generalizes: radiation is a boundary phenomenon in the continuum, the finite result of two chapters ago
lifted to a theorem about the completed space.

That all of this lifts intact is the payoff of having built the completion honestly. Because the
completion was assembled from the finite structure rather than imported over it, the finite theorems do
not have to be re-proved in a foreign setting; they rise with the space, each finite fact tanged up to
its continuum counterpart by the same density that makes the skeleton dense. The experiment on the
continuum limit names this passage: the finite model's features survive the limit when the limit is taken
as a genuine closure of the finite data, and here they demonstrably do. The construction does not arrive
at the continuum and hope its hard-won finite facts still hold; it carries them across the threshold,
proved.

The mechanism that carries each finite fact across is always the same, and it is worth naming because it
is why the lift is automatic rather than lucky. A property proved on the finite skeleton, if it is the
kind that respects the norm --- if small changes in the inputs make small changes in the property --- is
determined on the whole completion by its values on the dense skeleton, because every completion point is
the limit of skeleton points and the property follows them to the limit. Orthogonality respects the norm,
so it lifts. The vanishing of the interior mixed term respects the norm, so it lifts. The boundary trace
is built to be continuous, so it agrees with the limit of the finite traces. Each lift is an instance of
one principle: a continuous fact on a dense subset is a fact on the whole. The construction does not get
lucky chapter by chapter; it gets, once, a completion in which its finite facts are continuous, and
continuity does the rest. This is also why the radiation theorem generalizes rather than merely
surviving --- the continuum statement is the finite one's continuous extension, true everywhere because
true densely and continuous between.

\section{What the door grants, and what it does not}
\label{se:what-the-door-grants}

It is worth being precise about the one thing the open door still rests on, now that it is clear how
little that is. The construction builds the norm, the skeleton, the Cauchy sequences, the embedding, and
the lifts; all of these are assembled, not assumed. What it grants --- the only thing --- is that the
Cauchy sequences reach their limits, that the points of the completion exist as the generic the countable
skeleton approaches. This is not a new assumption dragged in at the door; it is the same standing grant
the book has carried since the tenth chapter, the genericity that the approached is reached, now wearing
its completion-shaped coat. The door does not import analysis; it asks only what the construction has
always asked, that the finite cuts arrive at their generic. Naming the completion has, like naming the
forcing, added nothing --- it has shown that the one place the construction seemed to import a world was
only the one place it exercises its single grant.

\section{What is demonstrated, and what is assumed}
\label{se:demonstrated-assumed-ch24}

The accounting is the cleanest the completion could hope for, because the chapter's whole point is how
little crosses into the assumed column.

What is demonstrated is the completion and the lift. The Sobolev--Hilbert door is built from the energy
norm and the separable skeleton, with no analytic import; the finite residue is proved to map faithfully
into it. And every finite structural fact is proved to lift: the split's orthogonality, the absence of an
interior mixed term, the agreement of the boundary trace with the terminal residual, the Cauchy
convergence in the Sobolev norm, and the generalized continuum theorem that radiation lives at the
boundary. Each is a checked statement, finite below and lifted above.

\begin{quote}
\emph{A note on the one grant, at the door.} The honest claim of this chapter is a negative one worth
stating positively. The completion does \emph{not} import the analysis of the real line --- not its
construction of the continuum, not its limits taken as given, none of it. It is built from the
construction's own separable space, and the only thing it asks beyond what it builds is the standing
genericity: that the countable Cauchy sequences reach the limits they approach. That single grant is the
genericity of the forcing, named two chapters ago, and the completion is just that genericity made into a
space. So the door long carried as the construction's one import turns out to import nothing; it spends,
at the threshold, exactly the one grant the book has spent everywhere, and not a theorem of analysis
more.
\end{quote}

What is assumed, then, is the standing grant alone, now shown to be all the completion ever needed. With
the door open, the finite and the continuum are one space, and every result the construction proved on
the finite side stands on the complete side as well. The construction has reached the completed setting
it pointed at from its first refinement, carrying its physics and its bookkeeping across intact. The
remaining chapters step outside the construction to read it from without --- to ask what an observer who
did not build it could measure of it --- beginning with the reading the next chapter takes from the
outside.
